\chapter*{General Introduction}
\addcontentsline{toc}{chapter}{General Introduction}
\markboth{General Introduction}{}

\lettrine[]{\textit{I}}{n modern financial software}, data quality plays a decisive role in the reliability of testing activities. Trading and portfolio management platforms manipulate many interdependent entities, such as accounts, securities, orders, executions, prices, and positions. Preparing coherent datasets for such systems is therefore a demanding task: every generated record must respect technical database constraints as well as financial business rules.

Within Linedata, the Longview platform is used to support front-office investment management workflows. Its quality assurance process requires realistic and reproducible datasets in order to validate new features, reproduce bugs, and run regression scenarios. However, preparing these datasets manually through SQL scripts is time-consuming, error-prone, and difficult to reuse across testing cycles. These limitations reduce productivity and make it harder to cover a wide range of financial scenarios.

The objective of this final-year project is to design and implement a configurable Test Data Generator for Longview. The proposed system provides a web interface that allows QA engineers and developers to define generation parameters without writing SQL manually. It supports standalone generation, connected generation based on existing SQL Server data, validation of generated records, configuration persistence, and optional Large Language Model enrichment to improve the realism of generated financial values.

The project was carried out using an Agile Scrum approach, with incremental sprints focused on data model discovery, generation engine implementation, connected mode implementation, LLM integration, configuration library implementation , and validation. This structure allowed the solution to evolve progressively according to technical findings and feedback from the target users.

This report is structured as follows:
\begin{itemize}
    \item \textbf{Chapter 1} presents the project context, the host organization, the Longview platform, the current test data preparation process, and the problem statement.
    \item \textbf{Chapter 2} details the requirements analysis, the selected methodology, the global architecture, the product backlog, and the technical environment.
    \item \textbf{Chapter 3} describes Sprint 1, dedicated to Longview data model discovery and the mapping of the target entities.
    \item \textbf{Chapter 4} presents Sprint 2, which delivered the core generation engine and the standalone mode.
    \item \textbf{Chapter 5} presents Sprint 3, which added SQL Server connectivity and connected generation from grabbed Longview data.
    \item \textbf{Chapter 6} presents Sprint 4, dedicated to optional Large Language Model enrichment for more realistic generated values.
    \item \textbf{Chapter 7} presents Sprint 5, which finalized the configuration library and the end-to-end application workflows.
\end{itemize}

The following chapters therefore document both the functional need behind the project and the technical decisions made to build a domain-aware tool capable of generating coherent Longview test data.
